\chapter*{前言}
\addcontentsline{toc}{chapter}{前言}

微积分课程使人学会计算极限、导数、积分与级数，而数学分析进一步追问：这些运算为什么成立，它们依赖实数的哪些性质，一个结论的条件究竟能削弱到什么程度，以及当熟悉的结论失效时，障碍发生在哪里。本书写给已经系统学过工科高等数学或同等课程，希望从计算性的微积分进入严格分析理论的读者。它不把读者当作微积分初学者，也不假定读者已经接受过数学系证明训练；全书试图搭建的，正是这两种经验之间常被省略的桥梁。

本书的目标不是把高等数学重新讲一遍。对求导、积分、常微分方程等常规计算，只在它们能够揭示定义、证明结构或适用边界时才作讨论。叙述的重心放在四个方面：概念为何产生，定义为何采取当前形式，主要定理如何彼此依赖，以及缺少条件时结论如何失败。对于高等数学中常作为直觉或计算规则使用的内容，本书将区分其严格命题、证明所需条件和可推广的结构。

全书由经典数学分析逐步延伸到现代分析。前半部分从实数系统、极限、连续、微分、积分、级数和多元分析出发，重新建立传统微积分的逻辑基础；随后进入度量空间、函数空间、测度与积分、各种函数收敛方式、\(L^p\) 空间、Fourier 分析和泛函分析；最后说明这些思想怎样进入常微分方程、Sobolev 空间、弱解与分布理论。这样的安排并不是把若干高级名词依次堆叠，而是沿着几条反复出现的主线展开：完备性保证极限对象存在，紧致性把无限过程化为有限控制，连续性表达结构在极限下的稳定，线性化把局部非线性问题转化为可处理的线性问题，函数空间则使“函数列的极限”成为空间中的收敛问题。

本书采用“动机—定义—例子与反例—定理—证明—条件辨析—后续联系”的基本组织方式。定义不是术语解释，而是后续推理的起点；证明不是对结论的附注，而是说明结论为何成立、条件在何处被使用的主要载体。重要定理尽量给出完整证明；若某个证明必须依赖尚未建立的理论，会明确标出依赖并在适当位置补全。反例与正面例子具有同等地位：正面例子说明定义能够捕捉什么，反例则划定定理的边界，并纠正由有限维、光滑或计算性经验造成的误判。

习题是正文的一部分，而不是章节末尾的附属材料。每章习题按概念与基本证明、结构性问题、拓展与反思分组，并依据章节在全书中的重要性配置训练量。部分题目补充课堂叙述中略去的技术细节，部分题目要求迁移定理或构造反例，另有少量项目题通过分步问题引向经典结果。正文仅给出必要提示，完整解答收入独立的配套解答册。使用本书时，建议在阅读证明前先尝试找出结论所需的关键定义，在阅读后再逐行标明每个条件的使用位置；对于反例，应同时检查对象范围、全部假设和失败的结论。

严格性并不意味着在第一页列出全部形式逻辑和集合论公理。本书在第0章建立后文实际使用的逻辑、集合、映射与证明语言，并明确所采用的集合论基础；完整的公理表、选择公理及其等价形式集中放在附录中。这样做既避免基础问题被含混带过，也避免集合论形式体系遮蔽分析学本身。除非特别说明，全书在经典逻辑和通常的 Zermelo--Fraenkel 集合论加选择公理的背景下工作；某个结论若只需更弱的选择原则，或其证明具有明显的非构造性，将在相应位置说明。

本书预设读者熟悉一元与多元微积分的基本计算、数项级数、幂级数、多重积分以及常微分方程的常见求解方法。必要的线性代数与集合论内容会在正文或附录中补充。阅读时不必一次掌握所有拓展章节；主线章节负责建立后续不可缺少的概念与定理，拓展专题则用于展示理论边界和更高层次的统一观点。

本书仍处于持续编写与审校之中。章节之间的依赖、符号体系、习题难度和证明密度将随全书推进不断接受一致性检查。希望最终形成的不是一本只供查阅结论的手册，而是一条可以实际行走的路线：从会使用微积分，走向能够理解、证明并继续推广分析学中的基本结论。
